When strain stretches a vortex, incompressibility makes the material tube
narrow as it lengthens. If circulation is retained, the narrowing makes it
spin faster. Viscosity spreads rotation into the surrounding fluid, so the
tube can narrow while losing circulation. Its vorticity
\(\omega=\nabla\times u\) satisfies
\[
 \partial_t\omega+(u\cdot\nabla)\omega
 =(\omega\cdot\nabla)u+\nu\Delta\omega.
\]
The hope is simple: viscosity smooths the fluid faster than nonlinear
transport can sharpen it.

A smaller core gives viscosity less distance to act across, but the faster
flow can also sharpen it more quickly. Under the Navier--Stokes rescaling
with \(\lambda>1\),
\[
 \begin{aligned}
 u_\lambda(x,t)&=\lambda u(\lambda x,\lambda^2t),\\
 p_\lambda(x,t)&=\lambda^2p(\lambda x,\lambda^2t),
 \end{aligned}
\]
transport and diffusion scale together at fixed \(\nu\). Smaller scales give
viscosity no automatic advantage.

The critical Sobolev order is \(s=\tfrac12\):
\[
 \|u_\lambda(\cdot,t)\|_{\dot H^s}
 =\lambda^{s-1/2}\|u(\cdot,\lambda^2t)\|_{\dot H^s}.
\]
Concentration makes the norm grow above that order and shrink below it;
at the critical order it stays invariant. Energy lies below this threshold.
The flow can become faster and more concentrated while its total energy
falls:
\[
 E[u_\lambda](t)=\lambda^{-1}E[u](\lambda^2t),
 \qquad E=\tfrac12\|u\|_2^2.
\]
An energy bound can survive the concentration we need to exclude.

Both transport and diffusion act over successively shorter times. If the
flow could keep repeating this concentration, those times would permit a
Zeno cascade at a finite \(T_*\):
\[
 \tau_j=\lambda^{-2j}\tau_0,
 \qquad \sum_{j=0}^{\infty}\tau_j<\infty.
\]
Scaling allows that possibility; the regularity
problem is whether the fluid can actually carry it out.
